\documentclass[11pt,a4paper]{article}
\usepackage[T1]{fontenc}
\usepackage[utf8]{inputenc}
\usepackage[spanish,es-noquoting,es-noshorthands]{babel}
\usepackage{lmodern,amsmath,amssymb,geometry,hyperref}
\geometry{margin=28mm}
\setlength{\emergencystretch}{3em}
\newcommand{\gr}{\operatorname{gr}}
\newcommand{\Hom}{\operatorname{Hom}}
\newcommand{\Exalcotop}{\operatorname{Exalcotop}}
\title{Notas matemáticas sobre los preliminares de EGA IV-1}
\author{}
\date{}
\begin{document}
% BEGIN RIGHTS AND AI AUTHORSHIP NOTICE
\begingroup
\small
\noindent\textbf{Derechos, disponibilidad y autoría}

Los derechos sobre las obras originales y sobre las contribuciones protegidas a estas traducciones y transcripciones derivadas corresponden a sus respectivos titulares. En la medida en que el editor tenga derechos sobre la traducción, la transcripción, las correcciones y la composición producidas por IA, los dedica al dominio público bajo CC0 1.0 Universal (https://creativecommons.org/publicdomain/zero/1.0/deed.es); esto no afecta a los derechos sobre las obras originales ni concede permiso para usarlas. Las menciones de derechos de terceros y las autorizaciones ya concedidas siguen siendo aplicables.

Estas ediciones se ofrecen gratuitamente, con fines no comerciales y exclusivamente educativos, bajo la presunción de que los respectivos titulares de derechos no se opondrían a que estas traducciones y transcripciones estuvieran disponibles gratuitamente junto a los originales franceses en \href{https://www.numdam.org/}{NUMDAM}. Esta presunción no significa que se haya obtenido autorización; el acceso gratuito a los originales no autoriza por sí solo su traducción, adaptación o republicación.

Si usted es titular de derechos sobre algún material presentado aquí y se opone a su disponibilidad, puede \href{https://github.com/KokunoYumeto/ega-fga-es/issues}{abrir una incidencia en GitHub}. El responsable de la publicación retirará el material afectado a petición del titular de los derechos.

Estas notas matemáticas y su composición están producidas por IA. El trabajo documentado incluye OpenAI Codex, GPT-6 Astra, esfuerzo Ultra. Se conserva la autoría de las obras originales y de las contribuciones citadas; no se atribuyen a la IA. El responsable de la publicación no tradujo el texto.
\par
\endgroup
\clearpage
% END RIGHTS AND AI AUTHORSHIP NOTICE

\maketitle
Estas notas se refieren al capítulo~0 de A.~Grothendieck y J.~Dieudonné,
\emph{Éléments de géométrie algébrique IV, première partie},
\emph{Publications Mathématiques de l'IHÉS} \textbf{20} (1964).
La numeración de los resultados es la de esa obra. La traducción según el
impreso y la edición con enmiendas son textos distintos; las justificaciones
matemáticas de las lecturas de esta última se exponen aquí por separado.

\section*{Regularidad y homotecias: 15.1.8}
La regularidad de un elemento sobre un módulo significa que la homotecia
correspondiente es inyectiva. Una extensión polinómica, considerada como módulo
sobre el anillo de base, es una suma directa de copias del módulo inicial;
conserva, por tanto, esa inyectividad. No se requiere sobreyectividad.
Por ejemplo, sean $A=k[x,y]$ y $M=A$, con la filtración $M_0=M$ y
$M_i=0$ para $i\geq1$. La sucesión $(x,y)$ es regular, pero en
$M/xM=k[y]$ la multiplicación por $y$ no tiene a $1$ en su imagen.
En el argumento sobre $\gr_\bullet(M)\otimes_A(A/\mathfrak J'')$
interviene, pues, la inyectividad. La biyectividad de los homomorfismos
$\psi$, $\psi'$ y del homomorfismo de (15.1.8.1) es una afirmación distinta.

\section*{Dimensión de un cociente: 16.3.4}
Bajo las hipótesis de (16.3.4), sea $N=M/xM$. El número de elementos
elegidos en el radical debe realizar el mínimo $s(N)$ para el que el cociente
de $N$ tiene longitud finita. Por (16.2.3),
$n=s(N)=\dim(N)$. Si esos elementos son $x_1,\ldots,x_n$, el cociente
\[
M/(xM+x_1M+\cdots+x_nM)
\]
tiene longitud finita. En consecuencia, $s(M)\leq n+1$ y
$\dim(M)\leq\dim(N)+1$. Una elección arbitraria, no mínima, de $n$
no bastaría por sí sola para obtener la última desigualdad.

\section*{Cocadenas de Hochschild: 18.4.4}
La estructura de $B$-módulo de $B\otimes_A B$ se define mediante el primer
factor. Si $h':B\otimes_A B\to L$ es $B$-lineal, la aplicación
$g(y)=h'(1\otimes y)$ es $A$-lineal. Recíprocamente, toda aplicación
$A$-lineal $g:B\to L$ define un homomorfismo $B$-lineal mediante
$h'(x\otimes y)=xg(y)$, pues $g(ay)=ag(y)$ para $a\in A$.
Esto no implica que $g$ sea $B$-lineal. Para $A=\mathbf Q$ y
$B=L=\mathbf Q[t]$, la derivación $g=d/dt$ satisface $g(1)=0$ y $g(t)=1$.
La fórmula del coborde es
\[
h'(d_2(1\otimes x\otimes y))=xg(y)-g(xy)+yg(x),
\]
con $g$ simplemente $A$-lineal.

\section*{Ideales abiertos de los dos anillos: 18.5.1}
Para el homomorfismo continuo $\rho:A\to B$, los cocientes considerados
son $A/\mathfrak J$ y $B/\mathfrak R$, donde
$\mathfrak J\subset A$, $\mathfrak R\subset B$ y
$\rho(\mathfrak J)\subset\mathfrak R$. Los refinamientos satisfacen
$\mathfrak J'\subset\mathfrak J\subset A$ y
$\mathfrak R'\subset\mathfrak R\subset B$.
En el orden $(\mathfrak R',\mathfrak J')$, los anillos ambientes son,
respectivamente, $B$ y $A$. Estas inclusiones proporcionan los homomorfismos
entre cocientes del sistema inductivo de (18.5.1.1).

\section*{Reducción y levantamiento de una base: 19.1.10}
La reducción de $u:M\to N$ módulo el ideal maximal $\mathfrak J_0$
tiene como módulo de llegada $N/\mathfrak J_0N$.
En la parte local de la demostración, sea
$r=\dim_{A/\mathfrak J}(M/\mathfrak JM)$ y elijamos $L=A^r$.
Los levantamientos de una base de $M/\mathfrak JM$ definen $f:L\to M$,
cuya reducción es un isomorfismo. Así, $L$ es libre de tipo finito.
La libertad de $L$ es la propiedad utilizada en el argumento del determinante
para la reducción de $u\circ f$; no basta con que $L$ sea de tipo finito.

\section*{Denominadores de un inverso local: 19.1.12}
Tras la reducción al caso en que $N$ es libre de tipo finito, un homomorfismo
$w_{\mathfrak p}:N_{\mathfrak p}\to M_{\mathfrak p}$ se escribe
\[
w_{\mathfrak p}=w\otimes(1/f)1_{A_{\mathfrak p}},
\qquad w:N\to M,\quad f\notin\mathfrak p.
\]
Si $w_{\mathfrak p}u_{\mathfrak p}=1$, entonces
$(wu)_{\mathfrak p}=f1_{M_{\mathfrak p}}$.
Como $M$ es de tipo finito, existe un mismo $g\notin\mathfrak p$ que
anula $wu-f1_M$ sobre todos sus generadores. Después de localizar en $fg$,
la multiplicación por $f$ es invertible y $f^{-1}w$ proporciona el inverso
por la izquierda. Es el homomorfismo local $w_{\mathfrak p}$ el que se
representa mediante la fracción indicada.

\section*{Prolongamiento al completado: 19.3.6}
El prolongamiento $B\to\widehat B\to C$ utiliza un homomorfismo
\emph{continuo} de $B$ al álgebra discreta $C$.
Un grupo topológico discreto es separado y completo; la continuidad da el
prolongamiento al completado separado y la densidad de la imagen de $B$
asegura su unicidad. La multiplicatividad y la compatibilidad con la acción
de $A$ se prolongan junto con las operaciones continuas.
Esta propiedad universal no se aplica a homomorfismos abstractos sin
hipótesis de continuidad.

\section*{Módulo de una extensión infinitesimal: 19.4.4}
Si $\mathfrak J^2=0$ en $C$, el producto fibrado
\[
E_\lambda=B_\lambda\mathbin{\times}_{C/\mathfrak J}C
\]
es una extensión de $B_\lambda$ por $\mathfrak J$.
La acción de $B$ sobre su núcleo pasa por
$B\to B_\lambda\to C/\mathfrak J$. En particular, ese núcleo está
anulado por el ideal abierto que define $B_\lambda$.
El grupo cuya anulación trivializa la extensión después de refinar el índice
es, por tanto, $\Exalcotop_A(B,\mathfrak J)$.
La hipótesis de (19.4.4) se aplica precisamente a este módulo discreto.

\section*{Anillo de coeficientes: 19.8.8}
La demostración utiliza el homomorfismo $u:W\to A$ elegido al comienzo;
no necesita una afirmación adicional de unicidad.
El homomorfismo local $w:B\to A$, con
$B=W[[T_1,\ldots,T_{m-1}]]$, hace de $A$ un $B$-módulo cuasifinito.
La completitud de $B$ y la separación de $A$ permiten aplicar
$(0_{\mathrm I},7.4.1)$ para concluir que $A$ es un $B$-módulo de tipo
finito. El anillo de escalares de esa conclusión es $B$.

\section*{Cocientes formalmente lisos: 19.10.3(i)}
Sea $\rho:A\twoheadrightarrow B=A/I$ un homomorfismo sobreyectivo.
En el caso discreto, $B$ es formalmente liso sobre $A$ si y solo si $I=I^2$.
En efecto, la aplicación idéntica de $B$ a $A/I$ debe levantarse a
$A/I^2$ como homomorfismo de $A$-álgebras. Esto obliga a que la imagen de
$I$ en $A/I^2$ sea nula, de donde $I=I^2$.
Recíprocamente, en un problema de levantamiento con ideal de cuadrado nulo
$J\subset C$, la imagen de $I$ en $C$ está contenida en $J$.
Si $I=I^2$, esa imagen está contenida en $J^2=0$; el homomorfismo
$A\to C$ se factoriza, pues, por $B$ y proporciona el levantamiento.
No es necesario que $\rho$ sea biyectivo.

\section*{Homomorfismos de extensiones: 20.6.25}
En la sucesión (20.6.25.1), el homomorfismo $u^1$ olvida parte de la
estructura de álgebra: una extensión relativa a $B$ se considera relativa
a $A$. Es la construcción de (18.3.6.4). El homomorfismo $v^1$ toma la
imagen recíproca de una extensión mediante $B\to C$, según (18.3.6.2).
La propia proposición (20.6.25) remite a estos dos apartados.
El último párrafo de su demostración, en la página impresa 146, lleva
«(18.3.4.2)»; la referencia pertinente para los dos homomorfismos es
(18.3.6.4) y (18.3.6.2), respectivamente. La traducción según el impreso
conserva el localizador impreso.

\section*{Álgebra de base de la lisura formal: 22.1.1}
En la torre $\Lambda\to A\to B$, las condiciones relativas a $A$
caracterizan la lisura formal de $B$ como $A$-álgebra.
En el sentido recíproco, la lisura formal de $B$ sobre $A$ y la de $A$
sobre $\Lambda$ implican la de $B$ sobre $\Lambda$, por composición
de los levantamientos. Son estas dos estructuras relativas las que
intervienen en (19.3.5(ii)) y en la demostración de (22.1.1).

\section*{Cancelación y número de variables: 23.1.3.1 y 23.1.4}
En (23.1.3.1), la cancelación de $x$, que no es divisor de cero, conserva
el factor $s'$ de la igualdad precedente: el miembro derecho es
$x^{n-1}as'$.
En (23.1.4), la inducción para el anillo de series formales
$A[[T_1,\ldots,T_r]]$ recae sobre el número $r$ de variables.

\section*{Intersecciones de anillos de Krull: 23.2.3--23.2.8}
El homomorfismo inyectivo del lema (23.2.3) es
\[
(A'\otimes_A B)_{\mathrm{red}}
\longrightarrow (K'\otimes_A R)_{\mathrm{red}}.
\]
Su numeración dentro del lema es (23.2.3.1); el impreso lleva (23.1.8.1).
En (23.2.7.1), cada $\varphi_i^{-1}(B'_i)$ es un subanillo de $K$.
Junto con $A'=B'\cap K$, las condiciones de pertenencia dan
\[
A'=\bigcap_i\varphi_i^{-1}(B'_i).
\]
La intersección finita de estos anillos de Krull es, por tanto, un anillo
de Krull. En (23.2.8), el resultado utilizado para las clausuras integrales
es la proposición (23.2.7).
\end{document}
